% answers_short.tex —— 1977 年·江苏苏州地区初试卷 · 简短答案速查
% 由 tools/make_short_answers.py 自 answers.tex 的 \gksolhead/\gkanswer 提取，
% 并与 1977/1977_数学真题_江苏苏州地区初试_答案要点.md 逐题比对后生成，勿手改（改源 MD 后重跑）。

\begin{gknotice}
\textbf{使用说明}：本节只给最终答案便于快速核对，完整推导见解析版。答案为编写者独立推导，\textbf{非官方标准答案}。
\end{gknotice}

\gkpart{题目}
\begin{itemize}[leftmargin=2.6em]
  \item[\textbf{1.}] (1) $-3y-1$；(2) $(2x^2-y^2)(2x^2+y^2)$；(3) $a^2-b^2-c^2+2bc$；(4) $2c$；(5) $\sqrt2$
  \item[\textbf{2.}] (1) $x=-3$；(2) $x=1$ 或 $x=-2$；(3) $x=-1$，$y=2$；(4) $y\leqslant-2$
  \item[\textbf{3.}] (1) $\sqrt3-1$；(2) $\sin B=\dfrac{\sqrt5}{5}$，$\operatorname{tg}A=2$；(3) $5$ 平方厘米；(4) 见解析版；(5) $5\sqrt2$
  \item[\textbf{4.}] 月平均增长率为 $10\%$。
  \item[\textbf{5.}] (1) $S=\dfrac{\sqrt2}{4}x(10-x)$（$0<x<10$）；(2) $x=5$（$\mathrm{cm}$）时 $S$ 最大；(3) $S$ 的最大值为 $\dfrac{25\sqrt2}{4}$ 平方厘米
  \item[\textbf{6.}] 切线方程为 $x-\sqrt3y-1=0$ 和 $x+\sqrt3y-1=0$。
\end{itemize}
